
@@M0@@
et c'est un
isomorphisme sur chaque composante puisque @@M1@@ a été supposé strictement parfait, d'où le ii).

2.7. Sur @@M2@@, le foncteur qui à un @@M3@@-Module gradué associe sa @@M4@@-ième composante est un foncteur exact à valeurs dans @@M5@@; il définit donc un foncteur :
@@M6@@
On peut, pour tout @@M7@@, considérer le foncteur @@M8@@ défini par la tensorisation par @@M9@@. L'homomorphisme (1.2.1) définit un morphisme de foncteurs sur @@M10@@, d'où un morphisme pour les foncteurs correspondants entre catégories dérivées :
@@M11@@
Soit @@M12@@ un complexe @@M13@@-parfait. Nous allons montrer qu'il existe une suite de triangles distingués :
@@M14@@
où @@M15@@ est acyclique, et les @@M16@@ du type @@M17@@, les @@M18@@ étant des complexes parfaits au sens ordinaire sur @@M19@@.

Puisque @@M20@@ est @@M21@@-parfait, et que @@M22@@ est quasi-compact, il existe un recouvrement fini de @@M23@@, et un entier @@M24@@ tel que sur tout ouvert de ce recouvrement, @@M25@@ soit isomorphe à un complexe borné à termes localement libres de type fini sur @@M26@@, soit @@M27@@, tel que @@M28@@ pour @@M29@@.

